\documentclass[10pt,a4paper]{amsart}
\usepackage[margin=19mm]{geometry}
\usepackage{amsmath,amssymb,mathtools}
\usepackage[scr=rsfs,cal=euler]{mathalfa}
\usepackage{xfrac}
\usepackage[unicode,hidelinks,bookmarks=false]{hyperref}
\newcommand{\ie}{i.e.\ }
\newcommand{\cf}{cf.\ }
\newcommand{\resp}{respectively\ }
\newcommand{\Subsection}{\subsection}
\providecommand{\nobreakdash}{\nobreak\hbox{-}}
\setlength{\emergencystretch}{2em}
\multlinegap0pt
\usepackage{bm}
\usepackage{bbm}
\usepackage{dsfont}
\usepackage{xspace}
\usepackage{cancel}
\usepackage{mathtools}
\usepackage{amsthm}
\makeatletter
\@ifundefined{theorem}{\newtheorem{theorem}{Theorem}}{}
\@ifundefined{lemma}{\newtheorem{lemma}[theorem]{Lemma}}{}
\makeatother
\newcommand{\F}{\mathbb{F}}
\title[Linear independence of powers for polynomials]{Linear independence of powers for polynomials}
\author{Alexandru Cr\u{a}ciun}
\date{Source: arXiv:2507.10163v1 (CC BY 4.0)}
\begin{document}
\maketitle

\noindent\emph{Source and licence.} Linear independence of powers for polynomials. Version arXiv:2507.10163v1 (CC BY 4.0), distributed under the Creative Commons Attribution 4.0 International licence, as recorded in the arXiv licence metadata for this exact version and shown on the abstract page https://arxiv.org/abs/2507.10163v1. Full rights notice: licenses/arxiv-2507.10163-CC-BY-4.0.txt.\par\smallskip

\noindent\textit{Editorial scope.} Counted unit: the Lemma whose source label is \texttt{lem:projection} in the article (source0.tex lines 262--270, source label \texttt{lem:projection}), delivered together with the article's own complete original proof of it (source0.tex lines 272--339, one \texttt{proof} environment of 68 lines). No mathematics is written, repaired or re-derived: statement and proof are the article's own text line for line. Required context retained uncounted: none. Every definition, display and label of those lines is retained unchanged. Omitted from this package and not redistributed: 1--261, 271--271, 340--443. Nothing inside a retained excerpt is omitted or altered: every retained body is delivered exactly as the article's lines are written, so no omission receipt of any kind accompanies this package. Citations: the retained excerpts carry no citation command at all, so no bibliography entry of the article is transcribed and no reference is redistributed. Reference adaptation: none was needed and none was made; every reference inside the delivered unit resolves inside it, so no unresolved reference or citation is produced. Printed slips kept literal: no printed word, symbol, hypothesis or number of the counted statement or of its proof was altered. Layout and class: the article's own class is \texttt{amsart} with packages including amssymb, amsrefs; the wrapper is the lane's pinned \texttt{amsart} editorial wrapper at 10pt on A4, which restates the theorem-like environments the retained text needs and adds the article's own macro definitions that the retained text needs (\texttt{\textbackslash F}), with extra wrapper packages (bm, bbm, dsfont, xspace, cancel, mathtools). The delivered numbering is the unit's own. Assets: no retained span contains a figure environment, a graphics-inclusion command, a PDF-page inclusion, a table or a TikZ picture; no image file is copied into this package, the package declares no asset dependency and no image file is redistributed with it. Licence: version arXiv:2507.10163v1 of this article is distributed under the Creative Commons Attribution 4.0 International licence, as recorded in the arXiv licence metadata for this exact version and shown on the abstract page https://arxiv.org/abs/2507.10163v1. A full rights notice is delivered at licenses/arxiv-2507.10163-CC-BY-4.0.txt. Characters: every retained excerpt is pure ASCII, so the delivered engine, XeLaTeX with its default Latin Modern text font, reports no missing character.
\begin{lemma}
	\label{lem:projection}
	Let $ d,k>0 $ be positive integers. Given any $ p_1, \ldots, p_k
	\in \F[x_1, \ldots, x_d]$ pairwise linearly independent with $ \deg_{x_1}
	p_i > 0 $ for all $ i \in \{1, \ldots, k\} $,
	there exists a point $ (\alpha_2, \ldots, \alpha_d) \in \F^{d-1} $ such that
	$ \bar{p}_1(x_1) = p_1(x_1, \alpha_2, \ldots, \alpha_d), \ldots,
	\bar{p}_k(x_1) = p_k(x_1, \alpha_2, \ldots, \alpha_d) $ are also pairwise linearly independent.
\end{lemma}

\begin{proof}
	We look at the set of points $ (\alpha_2, \ldots, \alpha_d) \in \F^{d-1} $ such
	that some pair $ \bar{p}_i, \bar{p}_j $ is linearly dependent, and we call it $ V
	$. We show that $ V $ is a variety of dimension $ \dim V < k-1 $. Since we
	assumed $ \F $ to be of characteristic $ 0 $, it is infinite, hence
	it cannot happen that $ V = \F^{d-1} $. Thus, we can always find our desired
	point.

	Take any two polynomials $ p_i, p_j $ with $ i\neq j $. We can then write them as
	\begin{align*}
		p_i(x_1, \ldots, x_d) &= p_{i,m}(x_2, \ldots, x_d)x_1^m + \ldots +
		p_{i,0}(x_2, \ldots, x_d),\\
		p_j(x_1, \ldots, x_d) &= p_{j,n}(x_2, \ldots, x_d)x_1^n + \ldots +
		p_{j,0}(x_2, \ldots, x_d),
	\end{align*}
	and we know that $ p_{i,m} $ and $ p_{j,n} $ are not zero. Let $ (\alpha_2,
	\ldots, \alpha_k) \in \F^{k-1}$ and consider three cases:	
	\begin{enumerate}
		\item
			$ \bar{p}_i = p_i $ and $ \bar{p}_j = p_j $, whence they are linearly
		independent by assumption.
		 \item
			 $ \bar{p}_i = p_i $ but $ \bar{p}_j \neq p_j $. Then $ p_{i,l} $ are
			 constants for all $ l \in \{1, \ldots, m\} $. Assuming $ \bar{p}_i
			 $ and $ \bar{p}_j $ are linearly dependent we have that either:
			 \begin{enumerate}
			 	\item
					 $ \bar{p}_j = 0 $ in which case $ (\alpha_2, \ldots, \alpha_d) $ must
					 be a solution to the linear system of equations $ p_{j, l} = 0 $ for $ l
					 \in \{1, \ldots, n\} $, i.e.\ $ (\alpha_2, \ldots, \alpha_d) \in
					 Z(p_{j,1}, \ldots, p_{j, n}) $. The ideal $ (p_{j,1}, \ldots, p_{j,
					 n}) $ cannot be empty since at least one $ p_{j,l} $ is
					 non-constant (we assumed $ \bar{p}_j \neq p_j $).

				\item
					$ \bar{p}_j \neq 0 $ in which case $ \bar{p}_j = \beta
					\bar{p}_i $ for some non-zero $ \beta \in \F $. Then $
					(\alpha_2, \ldots, \alpha_d) \in Z(p_{j, 1} - \beta p_{i,
					1}, \ldots, p_{j, n} - \beta p_{i, n})$, which is non-empty.
			 \end{enumerate} 

		\item
			$ \bar{p}_i \neq p_i $ and $ \bar{p}_j \neq p_j	$. We have another
			two cases:
			\begin{enumerate}
				\item
					$ \bar{p}_i = 0 $ or $ \bar{p}_j = 0 $. In this case $
					(\alpha_2, \ldots, \alpha_d) \in Z(p_{i,1}, \ldots, p_{i,
					m}) \cup Z(p_{j,1}, \ldots, p_{j, n})$.
				\item 
					$ \bar{p}_j, \bar{p}_i \neq 0$. If they are linearly
					dependent, we have $ \bar{p}_j = \beta \bar{p}_i $ for some
					non-zero $ \beta \in \F $. This means that $ (\alpha_2,
					\ldots, \alpha_d) \in Z(p_{j,1} -\beta p_{i,1}, \ldots,
					p_{j,\max(n,m)} -\beta p_{i,\max(n,m)}) $.
			\end{enumerate}
	\end{enumerate}

	In either case, we see that if $ \bar{p}_i $ and $ \bar{p}_j $ are linearly
	dependent, then $ (\alpha_2, \ldots, \alpha_d) $ must lie in a variety of
	dimension strictly smaller than $ d-1 $, which we call $ V_{i,j} $. 

	Proceeding similarly for every pair of polynomials, we see that some of the
	resulting polynomials will be pairwise linearly dependent only if $
	(\alpha_2, \ldots, \alpha_d) \in \bigcup_{i < j} V_{i,j}  = V$. $ V $ is the
	finite union of varieties whose dimension is smaller than $ d-1 $, hence it
	is also a variety of dimensions strictly smaller than $ d-1 $. 
\end{proof}

\end{document}
